On the domain $\Omega = \{(x,y)\mid x > 1\}$, the linear partial
differential operator
\begin{equation}
L
=
\frac{\partial^2}{\partial x^2}
+
\biggl(x-\frac1x\biggr)\frac{\partial^2}{\partial x\,\partial y}
-
\frac{\partial^2}{\partial y^2}
\label{90000033:operator}
\end{equation}
of second order is given.
\begin{teilaufgaben}
\item
Is the operator $L$ elliptic, parabolic or hyperbolic?
\item
Which part of the boundary can influence the value in the point
$P(2,1)$, or in other words, assuming that the $x$-direction is the
timelike direction, which part of the boundary is in the past
of the point $P$?
\end{teilaufgaben}

\begin{loesung}
\begin{figure}[h]
\definecolor{darkred}{rgb}{0.8,0,0}
\definecolor{darkgreen}{rgb}{0,0.6,0}
\centering
\begin{tikzpicture}[>=latex,thick,scale=4]
\fill[color=blue!20]
	plot[domain=1:2,samples=40]
		({\x},{0.5*\x*\x-1})
	--
	plot[domain=0:1,samples=50]
		({2-\x},{-ln(2-\x)+1+ln(2)})
	--
	cycle;
\draw[->] (-0.1,0) -- (2.5,0) coordinate[label={$x$}];
\draw[->] (0,-0.7) -- (0,2) coordinate[label={right:$y$}];
\draw[color=gray,line width=1.2pt] (1,-0.7) -- (1,1.9);
\begin{scope}
\clip (1,-0.7) rectangle (2.4,1.9);
\foreach \c in {-0.6,-0.4,...,3.4}{
	\draw[color=darkred,line width=1.2pt]
		plot[domain=1:2.5,samples=20] ({\x},{-ln(\x)+\c});
}
\foreach \d in {-4,-3.8,...,2}{
	\draw[color=darkgreen,line width=1.2pt]
		plot[domain=1:2.5,samples=20] ({\x},{0.5*\x^2+\d});
}
\end{scope}
\draw[color=orange,line width=1.4pt] (1,-0.5) -- (1,{1+ln(2)});
\fill[color=blue] (2,1) circle[radius=0.02];
\draw (-0.02,1) -- (0.02,1);
\node at (0,1) [left] {$1$};
\node at (0,0) [below left] {$0$};
\node at (1,0) [below left] {$1$};
\draw (2,-0.02) -- (2,0.02);
\node at (2,0) [below] {$2$};
\node[color=blue] at (1.98,1) [right] {$P$};
\end{tikzpicture}
\caption{Characteristics of the operator~\eqref{90000033:operator}
\label{90000033:fig}}
\end{figure}
\begin{teilaufgaben}
\item
The symbol matrix is
\[
A
=
\begin{pmatrix}
    1     & \frac12(x-\frac1x) \\
\frac12(x-\frac1x) &   -1
\end{pmatrix}
\]
and has determinante
\[
\det A
=
-1 -\frac14\biggl(x-\frac1x\biggr)^2 < 0,
\]
thus the operator is hyperbolic.
\item
The equation for the characteristic curves can be derived from the 
symbol matrix as
\[
\dot{y}^2 - \biggl(x-\frac1x\biggr)\dot{y}\dot{x} - \dot{x}^2 = 0.
\]
Dividing by $\dot{x}$ and using $y'=\dot{y}/\dot{x}$ gives
\[
y^{\prime 2} + \biggl(\frac1x-x\biggr)y' - 1 = 0.
\]
This differential equation has the factorization
\[
y^{\prime 2}
+
\biggl(\frac{1}x-x\biggr)y'
-
1
=
\biggl(y'+\frac1x\biggr)(y'-x)=0.
\]
This means that the characteristic curves satisfy the ordinary 
differential equations
\begin{align*}
y'&=-\frac1x
&
y'&=x
\intertext{which have the solutions}
y&=-\log x + C
&
y&=\frac12x^2 + D.
\end{align*}
The characteristic curves are shown in figure~\ref{90000033:fig}.
The curves through the point $P=(2,1)$ have the parameters
\begin{align*}
1=-\log2+C &\Rightarrow C=1+\log 2
&
1=\frac122^2+D
&\Rightarrow
D=-1.
\intertext{The two characteristic curves cross the boundary $x=1$ at}
y&=-\log 1+1+\log 2
=
1+\log 2
&
y&=\frac121^2-1
=
-\frac12.
\end{align*}
Thus only the boundary values in the points in the set
\(
\{(1,y) \mid -\frac12\le y\le 1+\log 2\}
\),
shown in orange in figure~\ref{90000033:fig},
can influence the value of the solution in the point $P$.
\qedhere
\end{teilaufgaben}
\end{loesung}

\begin{bewertung}
Symbol matrix ({\bf A}) 1 point,
operator is hyperbolic ({\bf H}) 1 point,
differential equation of characteristics ({\bf D}) 1 point,
factorization or other method to convert them into two first order
differential equations ({\bf F}) 1 point,
characteristic curves ({\bf C}) 1 point,
domain of influence for the point $P$ ({\bf P}) 1 point.
\end{bewertung}

